\documentclass{report}
\usepackage{amsmath,amssymb}
\setlength{\parindent}{0mm}
\setlength{\parskip}{1em}
\begin{document}
\begin{center}
\rule{6in}{1pt} \
{\large Qing Chu \\
{\bf ITERATIVE WAVEFRONT RECONSTRUCTION IN ADAPTIVE OPTICS}}

400 Dowman Dr  \\ W401 \\ Atlanta \\ GA \\ 30322
\\
{\tt qchu@emory.edu}\\
James Nagy\\
Stuart Jefferies\end{center}

Obtaining high resolution images of space objects from
ground based telescopes is challenging, and often
requires computational post processing using image
deconvolution methods. Good reconstructions can be
obtained if the convolution kernel can be accurately
estimated. The convolution kernel is defined by the
wavefront of light, and how it is distorted as it propagates
through the atmosphere. In this paper we describe the wavefront
reconstruction problem, and more specifically, a new
linear least squares model that exploits information from
multiple measurements.


\end{document}
